\section{본\hspace{0.5em}론}

여기에 내용을 입력하세요. 장 번호는 로마자(Ⅰ, Ⅱ, Ⅲ, \ldots)로 표기한다
(투고규정 13조).

\subsection{각 절 제목}

여기에 내용을 입력하세요. 절 번호는 아라비아자(1, 2, 3, \ldots)로 표기한다.

\subsubsection{각 항 제목}

여기에 내용을 입력하세요.

\paragraph{각 항 제목}

여기에 내용을 입력하세요. 도서와 컨퍼런스 논문의 인용 예시는 다음과
같다\cite{craig2004robotics,nasridnov2013visual}. 학위논문과 웹사이트의 인용
예시는 다음과 같다\cite{shannon1938symbolic,wright2015website}.

수식은 다음과 같이 오른쪽에 번호를 붙인다.
\begin{equation}
y = ax + b.
\label{eq:example}
\end{equation}
식 (\ref{eq:example})과 같이 본문에서 수식 번호를 인용한다.

그림~\ref{fig:example}\,은 그림 배치 예시이다. 그림과 표의 내용(그림 안의
글자, 표 안의 내용)은 영문으로만 기재하며, 그림이나 글자의 크기는 읽기
불편하지 않을 수준이어야 한다 (투고규정 14조). 그림 제목은 그림 하단에 영문
제목(Fig.~n.)과 국문 제목(그림~n.)을 병기하고, 국문 제목 끝의 마침표는
삭제한다.

\begin{figure}[t]
\centering
\framebox[0.95\columnwidth][c]{\rule[-14mm]{0pt}{28mm}%
  \fontsize{9}{11}\selectfont Contents of the figure must be written in English}
\bicaption{Please put the title of figure here.}{국문 제목을 입력하세요}
\label{fig:example}
\end{figure}

표~\ref{tab:example}\,은 표 배치 예시이다. 표 제목은 표 상단에 영문
제목(Table~n.)과 국문 제목(표~n.)을 병기한다.

\begin{table}[t]
\centering
\bicaption{Please put the title of table here.}{국문 제목을 입력하세요}
\label{tab:example}
\begin{tabular}{lcc}
\toprule
Parameters & Value & Unit \\
\midrule
Modulation format   & 16-QAM & --  \\
Channel bandwidth   & 10     & MHz \\
LNA gain            & 20     & dB  \\
\bottomrule
\end{tabular}
\end{table}

여기에 내용을 입력하세요. 여기에 내용을 입력하세요. 여기에 내용을 입력하세요.
여기에 내용을 입력하세요. 여기에 내용을 입력하세요. 여기에 내용을 입력하세요.
